News coverage links energy technologies to costs, risks and promises \parencite[e.g.,][]{stephens2009windmedia,ali2025hydrogenframing}. The shares of such framings, or frame shares, describe coverage, not acceptance, but can inform firms, investors and agencies monitoring a technology's public debate. Policy needs such evidence, which is slow and costly to collect. Large language models (LLMs) could therefore help \parencite[pp.~1--2]{toetzke2023llmclimatetech}. Yet their frame shares are useful only if the instrument behind them is validated.

\paragraph{Motivation.}
Such validation is rarely complete. Many content-analysis validations omit human-label reliability \parencite{song2020validations}. Three computational framing studies translate multilingual material or divide it between models \parencite{toetzke2022energydiscourse,jumle2024findingframes,guo2023openframing}, and none validates one instrument in each language. Validations typically report accuracy or agreement \parencite[p.~2]{linder2026multicalibrated}. Yet labels used directly can substantially bias statistics computed from them, such as shares, even at 80 to 90\% accuracy \parencite{egami2023dsl}. No reviewed study validates one LLM instrument at frame-element level, per language and on aggregate shares (Chapter~\ref{sec:related_work}).

\paragraph{Research Questions.}
The guiding research question is how a multilingual pipeline based on an LLM can be designed and validated for extracting technology frames from international renewable energy news. A technology frame is a news frame attached to a named energy technology. It is measured through frame elements, concrete presence questions such as whether a paragraph links wind power to its landscape impact. The elements marked present give the frames, and each frame carries an evaluation (Chapter~\ref{sec:background}, Figure~\ref{fig:anatomy}). Four research questions (RQ1 to RQ4) operationalise the guiding question:
\begin{enumerate}
\item[RQ1] How should such a pipeline be designed, and why is each part the way it is?
\item[RQ2] How well do its labels agree with a human reference, per language and per label dimension?
\item[RQ3] How far do identical repeated executions of the pipeline differ from one another?
\item[RQ4] How far do the aggregate prevalences it would report differ from those of the human reference, and where does the sample size limit what can be reported?
\end{enumerate}

\paragraph{Objective and Approach.}
One frozen codebook of coding rules supplies definitions for Claude Sonnet 4.6's system prompt and output schema and for the coding sheets of a single human coder (the author). The pipeline submits each paragraph untranslated, in German or English, the two languages the coder reads. Evaluation uses a stratified probability sample of 120 paragraphs, 60 per language, design-weighted to the eligible frame. This sampling frame (not a news frame) holds 618 paragraphs from a purposive, keyword-prefiltered corpus (Figure~\ref{fig:overview}a). The first of three identical model executions is scored against the coder's first coding. With the first labels hidden, the coder re-coded 40 paragraphs before the model run, as registered. A supplementary re-code covered the other 80 after the results were known; analyses using it are exploratory. Together they form the second coding. \emph{Registered} refers to an internal pre-analysis specification recorded before evaluation coding and prediction, not a public registry (Appendix~\ref{app:protocol}).

% 2026-11-09, second slimming pass (G-70): caption articles and legend wording tightened; no key removed.
% Figure 1 (Introduction) -- the study in one picture. Panel a is drawn by hand in TikZ: the evaluation design, lean
% (restructuring plan v3, 2026-11-05, Fig 1a): the 618-paragraph eligible frame, the 120-paragraph probability sample,
% the model and the single coder applying the same frozen codebook, and the comparison; one dashed marker for the
% analyses added after the registered result (listed in Table tab:added_analyses, not here). The system architecture
% that implements this design is images/architecture.tex. Design counts only, no measured result; every count is
% checked against the records: 618 paragraphs and 120 sampled (tools/runs/tables/dsl/dsl_prevalence.json, N and
% sampled), 60 German and 60 English (sample manifest; tab_primary), three executions with the first scored
% (Section research design), 40 paragraphs re-coded blind, 20 per language (Section on the human reference).
% Panel b is generated by tools/pipeline/figures.py results from the registered metrics (images/fig_headline.tex);
% no number in panel b is typed by hand. Style from the thesis preamble (figure style block); no internal codes
% (Chapter 1). The two panels sit flush left in one column, so that each panel letter, the left edge of its picture,
% lines up with the other.
% Deliberate deviation from the figure style specification (2026-11-05): panel b shows no displacement row; the
% largest downward displacement (18 points) is stated in the Findings paragraph beside the figure, and the Chapter 1
% budget leaves no room for a fourth row with its caption text. The author decides whether to add it.
% Palette 2026-11-06: panel a carries the thesis colour convention (preamble palette block). The model box is outlined
% in figModel (blue, 0.8pt, white fill), the coder box stays figInk, the comparison box, which combines both, is
% outlined in figAdjusted (green, 0.8pt) over its grey fill, the "later analyses" box and its arrow are figChanged
% (vermillion), the arrow from the model into the comparison is figModel; the arrow from the coder stays figInk. The
% frame and sample boxes stay black and grey. Every colour keeps its non-colour cue: the direct label in each box, the
% heavier outline on the two coloured boxes, the grey fill on the comparison, the dashed line and italic text on the
% later-analyses marker, and the position of each arrow at its source box. No count, label or layout changed; the
% caption names no colour in the panel descriptions and closes with the thesis-wide colour key.
\begin{figure}[htbp]
\centering
\begin{tabular}{@{}l@{}}
\begin{tikzpicture}[
  font=\figfontsmall,
  node distance=5mm and 7mm,
  pstep/.style={draw=figInk, line width=0.5pt, rectangle, rounded corners=1.5pt, minimum height=11mm,
    execute at begin node={\hyphenpenalty=10000\exhyphenpenalty=10000},
    align=center, inner sep=3pt, fill=white},
  key/.style={pstep, fill=figGreyC},
  modelbox/.style={pstep, draw=figModel, line width=0.8pt},
  cmpbox/.style={key, draw=figAdjusted, line width=0.8pt},
  later/.style={draw=figChanged, dashed, line width=0.5pt, rectangle, rounded corners=1.5pt, align=center,
    inner sep=3pt, text=figChanged, font=\figfontsmall\itshape,
    execute at begin node={\hyphenpenalty=10000\exhyphenpenalty=10000}},
  arr/.style={-{Latex[length=1.8mm]}, line width=0.6pt},
  darr/.style={arr, figChanged, dashed},
]
\node[pstep, text width=23mm] (frame) {Eligible frame\\618 paragraphs};
\node[key, text width=27mm, right=6mm of frame] (sample) {Probability sample\\120 paragraphs\\(60 German, 60 English)};
\node[modelbox, text width=36mm, above right=1.5mm and 10mm of sample.east] (model) {Model, codebook as prompt:\\three identical executions,\\the first scored};
\node[pstep, text width=36mm, below right=1.5mm and 10mm of sample.east] (coder) {Single coder, same codebook:\\reference labels;\\blind re-code of 40};
\node[cmpbox, text width=26mm, anchor=west] (cmp) at ([xshift=10mm]sample.east -| model.east) {Comparison:\\agreement, stability,\\displacement};
\draw[arr] (frame) -- (sample);
\draw[arr] (sample.east) -- ++(4mm,0) |- (model.west);
\draw[arr] (sample.east) -- ++(4mm,0) |- (coder.west);
\draw[arr, draw=figModel] (model.east) -- ++(4mm,0) |- ([yshift=1.5mm]cmp.west);
\draw[arr] (coder.east) -- ++(4mm,0) |- ([yshift=-1.5mm]cmp.west);
\node[later, text width=26mm, below=7mm of cmp] (later) {later analyses (Table~\ref*{tab:added_analyses})};
\draw[darr] (cmp) -- (later);
\node[panellabel] at ([xshift=-4mm, yshift=2mm]frame.west |- model.north) {a};
\end{tikzpicture}
\\[5mm]
% Generated by tools/pipeline/figures.py from a2_sonnet_20261013/metrics.json (sha256 9d69e490475a885f7216ceb5c49d295b9d7e55662b7092ca6faa1636e46c57f4). Do not edit by hand.
% Panel b of Figure 1 (paper/images/fig_overview.tex inputs this tikzpicture); element precision and recall from
% tools/runs/tables/model_comparison/comparison.json, the registered row bound to this metrics file by digest.
% Palette 2026-11-06: the precision and recall labels in model blue (figModel), as their open marks; the score marks take model blue from the preamble style pooled. Shapes and positions unchanged.
\begin{tikzpicture}
\begin{axis}[thesisfig,
  width=96mm, height=30mm, scale only axis,
  xmin=0, xmax=1, xtick={0,0.5,1}, xticklabels={0,0.5,1}, minor xtick={0.1,0.2,0.3,0.4,0.6,0.7,0.8,0.9}, xlabel={Score},
  ymin=0.4, ymax=3.80, y dir=reverse, rowaxis,
  ytick={1,2,3},
  yticklabels={{Technologies, micro-F1},{Evaluation in shared frames, accuracy},{Frame elements, micro-F1}},
]
\draw[figGreyB, line width=0.8pt] (axis cs:0.5572,3.42) -- (axis cs:0.7591,3.42);
\addplot[precisionmark] coordinates {(0.7591,3.42)};
\addplot[recallmark] coordinates {(0.5572,3.42)};
\addplot[pooled, error bars/.cd, x dir=both, x explicit] coordinates {(0.9350,1) -= (0.0336,0) += (0.0276,0) (0.8741,2) -= (0.0617,0) += (0.0498,0) (0.6427,3) -= (0.0622,0) += (0.0623,0)};
\node[figtext, right=3pt] at (axis cs:0.9626,1) {0.935};
\node[figtext, right=3pt] at (axis cs:0.9239,2) {0.874};
\node[figtext, right=3pt] at (axis cs:0.7050,3) {0.643};
\node[figtext, left=3pt, text=figModel] at (axis cs:0.5572,3.42) {recall 0.56};
\node[figtext, right=3pt, text=figModel] at (axis cs:0.7591,3.42) {precision 0.76};
\end{axis}
\node[panellabel] at ([yshift=2mm]current axis.north -| current axis.outer west) {b};
\end{tikzpicture}

\end{tabular}
% Caption revised 2026-11-06: takeaway restricted to the eligible, prefiltered frame (screening was not evaluated); panel b names its resampling unit; row labels with their metrics are requested in the generator of panel b.
\caption[Evaluation design and headline result]{\textbf{Against the first coding in prefiltered paragraphs, the model agrees closely on technologies but marks fewer frame elements.} \textbf{a},~Evaluation design. \textbf{b},~Registered agreement with the first coding, design-weighted, with article-resampled intervals; open marks: precision and recall, added later. Colours throughout: blue, registered model; black, coder; vermillion, changed or added later; green, model and coder combined; orange, model surplus.}
\label{fig:overview}
\end{figure}


\paragraph{Findings.}
Against the first coding, agreement is high for technology detection in prefiltered paragraphs (micro-F1 0.935) and for evaluations of shared frames (accuracy 0.874; Figure~\ref{fig:overview}b). Micro-F1 pools all labels and takes the harmonic mean of precision (the share of model marks the reference holds) and recall (the share of reference marks the model finds). End-to-end evaluation F1, also counting missed and invented frames, is 0.617. Repeated executions vary in output but give similar scores. Frame elements agree less, at a micro-F1 of 0.643 with no clear language difference, and the model marks fewer of them. Its prevalence, the share of paragraphs where it marks an element, lies below the first coding's in point estimate for 11 of the 15 reported elements, for investment by 18.3 percentage points.

Later analyses, added after the registered result was known (Table~\ref{tab:added_analyses}), set the recall deficit beside the coder's second coding. In point estimate, it narrows but persists on the 40 paragraphs re-coded before the model run. In exploratory comparisons with the full second coding, it nearly closes, but investment's share still differs by $-13.6$ [$-20.4$, $-7.6$] points. Neither coding is established as correct, so the central finding is a recall deficit against one coder, its size depending on the reference reading.

\paragraph{Contributions.}
The contribution is a tested combination of known components, not a new method. One frozen codebook is the common source of prompt, output schema and coding sheets, as proposed for codebooks that coders and models both read \parencite[p.~192]{halterman2026codebookllms}. A registered validation against one coder links label agreement per language and at element level, output variation across executions, the coder's stability over time and the displacement of aggregate shares. Later analyses indicate that the recall deficit's size, and the shares a design-based adjustment yields, depend on which coding is the reference.

\paragraph{Structure of the Thesis.}
Chapter~\ref{sec:background} defines technology frames and elements. Chapter~\ref{sec:related_work} reviews the validation gap. To close it, Chapter~\ref{sec:system_design} designs a pipeline and Chapter~\ref{sec:methodology} its evaluation. Chapter~\ref{sec:results} reports registered and later results, which Chapter~\ref{sec:discussion} interprets, before Chapter~\ref{sec:conclusion} answers RQ1 to RQ4.
